\chapter{Convolutional Neural Networks}

\section{Convolution as Structured Linear Operator}

Convolutional neural networks (CNNs) are designed to exploit structure in data, particularly in images and signals. The central operation is convolution, which can be understood as a structured linear transformation.

\medskip

\noindent
\textbf{Discrete convolution.}  
For a one-dimensional input $x \in \mathbb{R}^n$ and a kernel $w \in \mathbb{R}^k$, the convolution is defined as
\[
(y * w)_i = \sum_{j=1}^k w_j \, x_{i - j}.
\]

\medskip

\noindent
\textbf{Two-dimensional convolution.}  
For images, convolution is applied over spatial dimensions:
\[
y(i,j) = \sum_{u,v} w(u,v) \, x(i-u, j-v).
\]

\medskip

\noindent
\textbf{Linear operator view.}  
Convolution can be written as a matrix-vector multiplication:
\[
y = K x,
\]
where $K$ is a structured (Toeplitz or block-Toeplitz) matrix.

\medskip

\noindent
\textbf{Key property: parameter sharing.}
\begin{itemize}
    \item The same kernel is applied across all locations
    \item Reduces number of parameters
\end{itemize}

\medskip

\noindent
\textbf{Local connectivity.}
\begin{itemize}
    \item Each output depends only on a local region of the input
    \item Reflects spatial locality in data
\end{itemize}

\medskip

\noindent
\textbf{Insight.}  
Convolution imposes structure on linear transformations, making them efficient and suitable for structured data.

\section{Translation Invariance}

A key motivation for convolutional architectures is translation invariance.

\medskip

\noindent
\textbf{Translation equivariance.}  
A convolutional layer satisfies:
\[
x(\cdot - \tau) \;\rightarrow\; y(\cdot - \tau),
\]
meaning that shifting the input results in a corresponding shift in the output.

\medskip

\noindent
\textbf{Interpretation.}
\begin{itemize}
    \item Features detected at one location can also be detected elsewhere
    \item The model does not depend on absolute position
\end{itemize}

\medskip

\noindent
\textbf{From equivariance to invariance.}
\begin{itemize}
    \item Equivariance preserves spatial structure
    \item Invariance removes sensitivity to position
\end{itemize}

\medskip

\noindent
\textbf{Achieving invariance.}
\begin{itemize}
    \item Pooling operations
    \item Global averaging
\end{itemize}

\medskip

\noindent
\textbf{Insight.}  
Translation invariance reflects a key property of visual perception: objects can be recognized regardless of their location.

\section{Pooling and Hierarchies}

CNNs typically include pooling operations that reduce spatial resolution while preserving important features.

\medskip

\noindent
\textbf{Pooling.}
\begin{itemize}
    \item Max pooling: selects the maximum value in a region
    \item Average pooling: computes the average value
\end{itemize}

\medskip

\noindent
\textbf{Effect.}
\begin{itemize}
    \item Reduces dimensionality
    \item Introduces invariance to small transformations
\end{itemize}

\medskip

\noindent
\textbf{Hierarchical structure.}
\begin{itemize}
    \item Early layers detect simple features (edges, textures)
    \item Intermediate layers detect patterns (shapes, parts)
    \item Deeper layers detect complex structures (objects)
\end{itemize}

\medskip

\noindent
\textbf{Receptive field.}  
The region of the input that influences a given unit grows with depth.

\medskip

\noindent
\textbf{Insight.}  
CNNs build representations hierarchically, combining local features into global patterns.

\section{Applications in Vision}

CNNs have achieved remarkable success in visual tasks.

\medskip

\noindent
\textbf{Image classification.}
\begin{itemize}
    \item Assign a label to an image
    \item Example: object recognition
\end{itemize}

\medskip

\noindent
\textbf{Object detection.}
\begin{itemize}
    \item Identify and localize objects within an image
\end{itemize}

\medskip

\noindent
\textbf{Segmentation.}
\begin{itemize}
    \item Assign a label to each pixel
    \item Enables detailed understanding of scenes
\end{itemize}

\medskip

\noindent
\textbf{Generative modeling.}
\begin{itemize}
    \item Image synthesis
    \item Super-resolution
\end{itemize}

\medskip

\noindent
\textbf{Architectural examples.}
\begin{itemize}
    \item Convolutional layers + pooling
    \item Residual connections
    \item Encoder--decoder structures
\end{itemize}

\medskip

\noindent
\textbf{Insight.}  
CNNs are effective because they incorporate domain-specific structure into the model.

\section{A Unifying Perspective}

Convolutional neural networks combine several key ideas:

\begin{itemize}
    \item \textbf{Structured linear operators:} convolution
    \item \textbf{Inductive bias:} locality and translation invariance
    \item \textbf{Hierarchy:} multi-scale representation learning
\end{itemize}

\medskip

\noindent
\textbf{Core components.}
\begin{itemize}
    \item Convolutional layers
    \item Nonlinear activations
    \item Pooling operations
\end{itemize}

\medskip

\noindent
\textbf{Key insight.}  
By incorporating structure into the architecture, CNNs achieve both efficiency and strong generalization.

\section{Outlook}

This chapter introduced convolutional neural networks:
\begin{itemize}
    \item convolution as a structured linear operator,
    \item translation equivariance and invariance,
    \item hierarchical feature learning,
    \item applications in vision.
\end{itemize}

\medskip

In the next chapter, we study sequence models, which extend neural networks to sequential and temporal data.

\medskip

\noindent
\textbf{Guiding principle.}  
Effective models exploit the structure of the data they are designed to process.